\chapter{Modern Architecture Deep Dives}

\section{RMSNorm}

RMSNorm simplifies LayerNorm by removing mean subtraction and bias:

\begin{align*}
  \textbf{LayerNorm:} \quad & \hat{x}_i = \frac{x_i - \mu}{\sqrt{\sigma^2 + \epsilon}}, \quad y_i = \gamma \hat{x}_i + \beta \\[6pt]
  \textbf{RMSNorm:} \quad & \hat{x}_i = \frac{x_i}{\sqrt{\frac{1}{d}\sum_{j=1}^{d} x_j^2 + \epsilon}}, \quad y_i = \gamma \hat{x}_i
\end{align*}

Key differences:
\begin{itemize}
  \item RMSNorm removes \textbf{mean subtraction}: it only rescales, not
        recenters
  \item RMSNorm has \textbf{no bias} parameter ($\beta$), only a gain
        $\gamma$
  \item In practice, negligible speed difference---normalization is a
        tiny fraction of total compute
\end{itemize}

\begin{keyconcept}
  \textbf{Float32 upcasting}: Both LayerNorm and RMSNorm must upcast to
  float32 before computing mean/variance/RMS. Low-precision types (bf16,
  fp16) can overflow or lose precision in these reduction operations.
  After computing the normalization in float32, the result is cast back
  to the original dtype before the affine transform.
\end{keyconcept}

\section{RoPE (Rotary Positional Embeddings)}

RoPE encodes position by \emph{rotating} pairs of dimensions in the
query and key vectors, rather than adding positional embeddings.

\subsection{The Key Idea}

For each pair of dimensions $(2i, 2i+1)$ at position $m$:

\[
  \begin{bmatrix} q_{2i}' \\ q_{2i+1}' \end{bmatrix}
  =
  \begin{bmatrix} \cos(m\theta_i) & -\sin(m\theta_i) \\ \sin(m\theta_i) & \cos(m\theta_i) \end{bmatrix}
  \begin{bmatrix} q_{2i} \\ q_{2i+1} \end{bmatrix}
\]

where $\theta_i = \frac{1}{\text{base}^{2i/d}}$ (base = 10000 by
default).

This is a 2D rotation by angle $m\theta_i$.

\subsection{Efficient Implementation: Rotate-and-Interleave}

Instead of handling pairs explicitly, we use the identity:

\[
  q' = q \odot \cos(m\boldsymbol{\theta}) + \text{rotate\_half}(q) \odot \sin(m\boldsymbol{\theta})
\]

where $\text{rotate\_half}$ constructs $(-q_{2i+1}, q_{2i})$ for each
pair, then interleaves. In code:

\begin{pseudocode}[Apply Rotary Embedding]
Reshape $x$ to $[\ldots, d/2, 2]$ \hfill\textrm{\small (pair up dimensions)}\\
$x_{\text{rot}} \leftarrow$ stack $(-x[\ldots, 1], x[\ldots, 0])$ along last dim\\
$x_{\text{rot}} \leftarrow$ flatten $x_{\text{rot}}$ \hfill\textrm{\small (interleave back)}\\
\textbf{return} $x \odot \cos + x_{\text{rot}} \odot \sin$
\end{pseudocode}

\begin{keyconcept}
  \textbf{Critical property}: With RoPE, the dot product
  $q_m^\top k_n$ depends \emph{only on the relative distance} $(m - n)$,
  not on absolute positions. This is because the rotation angle for $q$
  at position $m$ minus the rotation angle for $k$ at position $n$ equals
  the angle for relative position $(m-n)$.

  \textbf{RoPE is applied only to Q and K, never to V}. Position
  information enters exclusively through attention scores.
\end{keyconcept}

\section{SwiGLU}

SwiGLU replaces the standard GELU FFN with a gated architecture:

\begin{align*}
  \textbf{GELU FFN:} & \quad \text{FFN}(x) = W_2\!\left(\text{GELU}(W_1 x)\right) \\[4pt]
  \textbf{SwiGLU:} & \quad \text{FFN}(x) = W_{\text{down}}\!\left(W_{\text{gate}}(x) \odot \text{SiLU}(W_{\text{up}}(x))\right)
\end{align*}

where $\text{SiLU}(x) = x \cdot \sigma(x)$ (also called ``swish'').

\begin{keyconcept}
  \textbf{Key differences from standard FFN}:
  \begin{itemize}
    \item 3 linear layers instead of 2 ($W_{\text{gate}}$,
          $W_{\text{up}}$, $W_{\text{down}}$)
    \item The gate controls which features pass through: when
          $\text{SiLU}(W_{\text{up}}(x)) \approx 0$, the gate ``closes''
          and blocks information
    \item Despite having more parameters per layer, SwiGLU is
          empirically more parameter-efficient---it matches or exceeds
          GELU FFN quality with fewer total parameters
  \end{itemize}

  LLaMA compensates for the extra projection by using
  $d_{\text{ff}} = \frac{2}{3} \cdot 4d_{\text{model}}$ instead of
  $4d_{\text{model}}$; our implementation keeps $4d_{\text{model}}$
  for simplicity.
\end{keyconcept}

\section{\texorpdfstring{GQA}{GQA} (Grouped Query Attention)}

GQA reduces the KV cache by sharing KV heads across groups of query
heads:

\begin{center}
\begin{tikzpicture}[
  node distance=0.3cm and 0.8cm,
  qnode/.style={rectangle, draw, fill=accent!30, minimum width=0.5cm, minimum height=0.5cm, font=\tiny},
  kvnode/.style={rectangle, draw, fill=highlight!30, minimum width=0.5cm, minimum height=0.5cm, font=\tiny},
  label/.style={font=\small\bfseries}
]
 
  \node[label] (mhlabel) {MHA};
  \foreach \i in {1,...,8} {
    \node[qnode, right=0.15cm of mhlabel] (q\i) at ({1.5+\i*0.4}, 0) {};
    \node[kvnode, below=0.2cm of q\i] {};
  }
  
  \node[label, right=1.2cm of q8] (gqlabel) {GQA};
  \foreach \i in {1,...,8} {
    \node[qnode] (gq\i) at ({7+\i*0.4}, 0) {};
  }
  \foreach \i in {1,2} { \node[kvnode] (gkv\i) at ({7.2+\i*1.6}, -0.7) {}; }
  \draw[-{Stealth}, gray] (gq1.south) -- (gkv1.north);
  \draw[-{Stealth}, gray] (gq2.south) -- (gkv1.north);
  \draw[-{Stealth}, gray] (gq3.south) -- (gkv1.north);
  \draw[-{Stealth}, gray] (gq4.south) -- (gkv1.north);
  \draw[-{Stealth}, gray] (gq5.south) -- (gkv2.north);
  \draw[-{Stealth}, gray] (gq6.south) -- (gkv2.north);
  \draw[-{Stealth}, gray] (gq7.south) -- (gkv2.north);
  \draw[-{Stealth}, gray] (gq8.south) -- (gkv2.north);
  
  \node[label, right=1.2cm of gq8] (mqlabel) {MQA};
  \foreach \i in {1,...,8} {
    \node[qnode] (mq\i) at ({12.5+\i*0.4}, 0) {};
  }
  \node[kvnode] (mkv) at ({14.1}, -0.7) {};
  \foreach \i in {1,...,8} {
    \draw[-{Stealth}, gray] (mq\i.south) -- (mkv.north);
  }
\end{tikzpicture}
\end{center}

\begin{center}
\begin{tabular}{lccc}
\toprule
 & \textbf{MHA} & \textbf{GQA} & \textbf{MQA} \\
\midrule
Query heads & $h$ & $h$ & $h$ \\
KV heads & $h$ & $h_{\text{kv}} < h$ & 1 \\
Group size & 1 & $h / h_{\text{kv}}$ & $h$ \\
KV cache savings & None & $\times\, h_{\text{kv}}/h$ & $\times\, 1/h$ \\
\bottomrule
\end{tabular}
\end{center}

\begin{keyconcept}
  \textbf{Implementation}: After computing K and V with the smaller
  $W_K$ and $W_V$ matrices ($d_{\text{model}} \rightarrow h_{\text{kv}}
  \cdot d_k$), we expand them using
  \texttt{repeat\_interleave} to match the number of query heads:

  \begin{enumerate}
    \item Compute K, V with reduced projection: $[b, s, h_{\text{kv}} \cdot d_k]$
    \item Reshape: $[b, s, h_{\text{kv}}, d_k]$
    \item Transpose: $[b, h_{\text{kv}}, s, d_k]$
    \item Expand: $[b, h, s, d_k]$ via \texttt{repeat\_interleave}
  \end{enumerate}

  RoPE is applied \emph{after} KV expansion to avoid subtle numerical
  issues. The special cases: $h_{\text{kv}} = h$ gives standard MHA,
  $h_{\text{kv}} = 1$ gives Multi-Query Attention (MQA).
\end{keyconcept}

\section*{Review Questions}

\begin{reviewquestion}[1]
  RMSNorm removes both mean subtraction and the bias parameter compared
  to LayerNorm. What is the intuition for why removing mean subtraction
  still works well?
\end{reviewquestion}

\begin{reviewquestion}[2]
  Prove that with RoPE, the dot product $q_m^\top k_n$ depends only on
  the relative position $(m - n)$, not on absolute positions $m$ and $n$
  individually. (Hint: use the 2D rotation formulation.)
\end{reviewquestion}

\begin{reviewquestion}[3]
  Consider a model with 32 query heads and 4 KV heads using GQA.
  During autoregressive generation, how much smaller is the KV cache
  compared to standard MHA? At what cost?
\end{reviewquestion}